% ============================================================
% AgroSmart Andino -- Presentacion Beamer
% Maestria en Inteligencia Artificial -- UNI
% Curso: Big Data | Grupo 2 | 15 minutos
% ============================================================
\documentclass[14pt,aspectratio=169]{beamer}

\usetheme{Madrid}
\usecolortheme{default}

\usepackage[utf8]{inputenc}
\usepackage[T1]{fontenc}
\usepackage[spanish,es-nodecimaldot,es-noindentfirst]{babel}
\usepackage{graphicx}
\usepackage{booktabs}
\usepackage{array}
\usepackage{amsmath}
\usepackage{amssymb}
\usepackage{xcolor}
\usepackage{colortbl}
\usepackage{listings}
\usepackage{tcolorbox}
\usepackage{multirow}

% Paleta de colores AgroSmart
\definecolor{agroverde}{RGB}{27,67,50}
\definecolor{agrolight}{RGB}{45,106,79}
\definecolor{agromid}{RGB}{82,183,136}
\definecolor{agroclear}{RGB}{184,233,199}
\definecolor{agroazul}{RGB}{58,134,255}
\definecolor{agrorojo}{RGB}{214,40,40}
\definecolor{agronaranja}{RGB}{244,162,97}
\definecolor{agrogris}{RGB}{108,117,125}

% Personalizacion del tema
\setbeamercolor{palette primary}{bg=agroverde, fg=white}
\setbeamercolor{palette secondary}{bg=agrolight, fg=white}
\setbeamercolor{palette tertiary}{bg=agromid, fg=white}
\setbeamercolor{palette quaternary}{bg=agroverde, fg=white}
\setbeamercolor{structure}{fg=agroverde}
\setbeamercolor{frametitle}{bg=agroverde, fg=white}
\setbeamercolor{title}{fg=white}
\setbeamercolor{block title}{bg=agrolight, fg=white}
\setbeamercolor{block body}{bg=agroclear!30}
\setbeamercolor{block title alerted}{bg=agrorojo, fg=white}
\setbeamercolor{block body alerted}{bg=agrorojo!10}
\setbeamercolor{block title example}{bg=agroazul!80, fg=white}
\setbeamercolor{block body example}{bg=agroazul!10}
\setbeamercolor{itemize item}{fg=agrolight}
\setbeamercolor{itemize subitem}{fg=agromid}
\setbeamercolor{enumerate item}{fg=agroverde}
\setbeamerfont{title}{size=\large, series=\bfseries}
\setbeamerfont{frametitle}{size=\normalsize, series=\bfseries}
\setbeamertemplate{navigation symbols}{}

% Footer
\setbeamertemplate{footline}{%
  \leavevmode%
  \hbox{%
    \begin{beamercolorbox}[wd=.40\paperwidth,ht=2.25ex,dp=1ex,center]{palette primary}%
      \usebeamerfont{author in head/foot}\insertshortauthor
    \end{beamercolorbox}%
    \begin{beamercolorbox}[wd=.35\paperwidth,ht=2.25ex,dp=1ex,center]{palette secondary}%
      \usebeamerfont{title in head/foot}\insertshorttitle
    \end{beamercolorbox}%
    \begin{beamercolorbox}[wd=.25\paperwidth,ht=2.25ex,dp=1ex,right]{palette tertiary}%
      \insertframenumber{}/\inserttotalframenumber\hspace*{2ex}
    \end{beamercolorbox}%
  }%
}

\lstset{
  basicstyle=\ttfamily\tiny,
  backgroundcolor=\color{gray!10},
  frame=single, framerule=0pt,
  breaklines=true, showstringspaces=false,
  keywordstyle=\color{agrolight}\bfseries,
  commentstyle=\color{agrogris},
  literate=
    {\'a}{{\'{a}}}1 {\'e}{{\'{e}}}1 {\'i}{{\'{i}}}1
    {\'o}{{\'{o}}}1 {\'u}{{\'{u}}}1 {\~n}{{\~{n}}}1
    {\'A}{{\'{A}}}1 {\'E}{{\'{E}}}1 {\'I}{{\'{I}}}1
    {\'O}{{\'{O}}}1 {\'U}{{\'{U}}}1
}

% Datos del documento
\title[AgroSmart Andino]{%
  \textbf{AgroSmart Andino}\\[0.2em]
  \small Pipeline Big Data con Dashboard en Tiempo Real\\
  para Monitoreo de Papa Nativa en Huangascar, Yauyos
}
\author[Grupo 2 -- Maestr\'ia IA UNI]{%
  Montes E. $\cdot$ Vizcardo F. $\cdot$ Massa Q.
  $\cdot$ Huali C. $\cdot$ Lazaro G.
}
\institute[UNI]{%
  Universidad Nacional de Ingenier\'ia\\
  Maestr\'ia en Inteligencia Artificial -- Big Data 2026\\
  \textit{Docente: Mg. Rosa Virginia Encinas Quille}
}
\date{Mayo 2026}

% ============================================================
\begin{document}
% ============================================================

\begin{frame}
  \titlepage
\end{frame}

\begin{frame}{Contenido}
  \tableofcontents[hideallsubsections]
\end{frame}

% ============================================================
\section{Introducci\'on y Problema}
% ============================================================

\begin{frame}{El problema: papa nativa bajo amenaza clim\'atica}
  \begin{columns}[T]
    \begin{column}{0.50\textwidth}
      \begin{block}{Contexto}
        \begin{itemize}
          \item Huangascar, Yauyos, Lima
          \item 3\,200 m.s.n.m. -- cultivo ancestral
          \item Papa nativa: variedad prioritaria
          \item Sin sistemas de alerta temprana
        \end{itemize}
      \end{block}
      \vspace{0.4em}
      \begin{alertblock}{Riesgos cr\'iticos}
        \begin{itemize}
          \item \textbf{Heladas:} 18 d\'ias cr\'iticos/a\~no
          \item \textbf{Tiz\'on tard\'io:} 47 d\'ias/a\~no
          \item \textbf{Estr\'es h\'idrico:} 31 d\'ias/a\~no
          \item P\'erdida 2022: \textbf{--38\,\%} de cosecha
        \end{itemize}
      \end{alertblock}
    \end{column}
    \begin{column}{0.48\textwidth}
      \begin{exampleblock}{Situaci\'on actual}
        \begin{itemize}
          \item Decisi\'on emp\'irica del agricultor
          \item Datos NASA POWER disponibles (gratis)
          \item AWS Academy con \$100 cr\'editos
          \item Oportunidad: Big Data + Cloud
        \end{itemize}
      \end{exampleblock}
      \vspace{0.4em}
      \begin{block}{Nuestra propuesta}
        \centering
        \textbf{Pipeline Big Data completo}\\
        con alerta temprana y recomendaci\'on\\
        de riego en tiempo cuasi-real
      \end{block}
    \end{column}
  \end{columns}
\end{frame}

% ── OBJETIVOS ───────────────────────────────────────────────
\begin{frame}{Objetivos del sistema AgroSmart Andino}
  \begin{block}{Objetivo general}
    Dise\~nar e implementar un \textbf{pipeline Big Data} para el
    monitoreo inteligente de papa nativa en Huangascar, con alerta
    temprana y recomendaci\'on de riego en tiempo cuasi-real.
  \end{block}
  \vspace{0.4em}
  \textbf{Objetivos espec\'ificos:}
  \begin{enumerate}
    \item Consolidar 3 fuentes heterog\'eneas en Data Lake S3
      \hfill{\footnotesize 72\,397 registros}
    \item Pipeline ETL distribuido con Apache Spark en EMR
      \hfill{\footnotesize 3\,653 integrados}
    \item Calcular 4 \'indices agroclim\'aticos para papa andina
      \hfill{\footnotesize IHB, ITT, IEH, ET$_0$}
    \item Modelos ML de riesgo fitosanitario y rendimiento
      \hfill{\footnotesize RF + GBT}
    \item Dashboard Streamlit con sem\'aforos y alertas en vivo
      \hfill{\footnotesize 11 KPIs}
  \end{enumerate}
\end{frame}

% ============================================================
\section{Arquitectura del Sistema}
% ============================================================

\begin{frame}{Arquitectura: flujo de 5 capas en AWS}
  \vspace{0.2em}
  \begin{center}
  \setlength{\tabcolsep}{3pt}
  \begin{tabular}{ccccc}
    \cellcolor{agroazul!80}\textcolor{white}{\footnotesize\textbf{Fuentes}} &
    \cellcolor{agronaranja!80}\textcolor{white}{\footnotesize\textbf{S3 Raw}} &
    \cellcolor{agrolight}\textcolor{white}{\footnotesize\textbf{EMR+Spark}} &
    \cellcolor{agroverde}\textcolor{white}{\footnotesize\textbf{S3 Gold}} &
    \cellcolor{agrogris}\textcolor{white}{\footnotesize\textbf{Cassandra}} \\
    \cellcolor{agroazul!80}\textcolor{white}{\tiny NASA/IoT} &
    \cellcolor{agronaranja!80}\textcolor{white}{\tiny 3 capas} &
    \cellcolor{agrolight}\textcolor{white}{\tiny ETL + ML} &
    \cellcolor{agroverde}\textcolor{white}{\tiny 4 CSV} &
    \cellcolor{agrogris}\textcolor{white}{\tiny 5 tablas} \\[0.3em]
    \multicolumn{5}{c}{$\Rightarrow\quad\Rightarrow\quad\Rightarrow\quad\Rightarrow$} \\
    \multicolumn{5}{c}{\cellcolor{agromid!40}
      \textbf{Streamlit Dashboard} | localhost:8501 | refresco 5 min}
  \end{tabular}
  \end{center}

  \vspace{0.3em}
  \begin{columns}[T]
    \begin{column}{0.32\textwidth}
      \begin{block}{Capa de ingesta}
        \footnotesize
        \textbf{S3 raw/processed/gold}\\
        Bucket: agrosmart-andino\\
        3 datasets, 6.3 MB
      \end{block}
    \end{column}
    \begin{column}{0.32\textwidth}
      \begin{block}{Capa de procesamiento}
        \footnotesize
        \textbf{EMR 6.15 -- Spark 3.4}\\
        1 master + 1 core\\
        m5.xlarge, JupyterHub
      \end{block}
    \end{column}
    \begin{column}{0.32\textwidth}
      \begin{block}{Capa de servicio}
        \footnotesize
        \textbf{Cassandra 4.1 (Docker)}\\
        5 tablas, TTL 90 d\'ias\\
        Streamlit dashboard
      \end{block}
    \end{column}
  \end{columns}
\end{frame}

% ── DATOS ───────────────────────────────────────────────────
\begin{frame}{Corpus de datos: 3 fuentes, 72\,397 registros}
  \begin{columns}[T]
    \begin{column}{0.55\textwidth}
      \begin{tabular}{lrrl}
        \toprule
        \textbf{Fuente} & \textbf{Filas} & \textbf{MB} & \textbf{Script} \\
        \midrule
        NASA POWER API & 18\,265 & 1.6 & 01\_download \\
        Sensores IoT sint. & 52\,632 & 4.6 & 02\_generate \\
        MIDAGRI hist\'orico & 1\,500 & 0.1 & 03\_prepare \\
        \midrule
        \rowcolor{agroclear!50}
        \textbf{Total validado} & \textbf{72\,397} & \textbf{6.3} & 04\_validate \\
        \bottomrule
      \end{tabular}

      \vspace{0.5em}
      \begin{block}{NASA POWER -- 7 variables diarias}
        \footnotesize
        T$_{\max}$, T$_{\min}$, T$_{\text{med}}$,
        Precipitaci\'on, HR, Radiaci\'on, Viento\\
        Rean\'alisis $0.5^{\circ}\times0.5^{\circ}$ -- 2015 al 2024
      \end{block}
    \end{column}
    \begin{column}{0.43\textwidth}
      \begin{exampleblock}{Cobertura geogr\'afica}
        \footnotesize
        \textbf{5 estaciones:}\\
        Huangascar (base, 3\,200m)\\
        Yauyos, Carania, Tanta, Vilca\\[0.3em]
        \textbf{3 nodos IoT:}\\
        Parcela Baja (3\,000 m)\\
        Parcela Media (3\,200 m)\\
        Parcela Alta (3\,500 m)
      \end{exampleblock}
      \vspace{0.3em}
      \begin{alertblock}{Validaci\'on 04}
        \footnotesize
        3/3 datasets aprobados\\
        Sem\'afico: \textbf{VERDE}\\
        0 duplicados, 0 nulos cr\'it.
      \end{alertblock}
    \end{column}
  \end{columns}
\end{frame}

% ============================================================
\section{Pipeline ETL con Apache Spark}
% ============================================================

\begin{frame}{Pipeline ETL: 3 fases en Spark 3.4 sobre EMR}
  \begin{columns}[T]
    \begin{column}{0.48\textwidth}
      \begin{block}{Fase 1 -- Ingesta y limpieza}
        \begin{itemize}\footnotesize
          \item Lectura 3 CSV desde S3 Raw
          \item \texttt{dropDuplicates()}
          \item Imputaci\'on: mediana/moda
          \item Tipado y fechas estandarizadas
        \end{itemize}
      \end{block}
      \vspace{0.25em}
      \begin{block}{Fase 2 -- Integraci\'on}
        \begin{itemize}\footnotesize
          \item Join clima + IoT por fecha/estaci\'on
          \item Agregaci\'on horaria a diaria
          \item 3\,653 registros integrados
        \end{itemize}
      \end{block}
      \vspace{0.25em}
      \begin{block}{Fase 3 -- Escritura Gold}
        \begin{itemize}\footnotesize
          \item Parquet particionado: \texttt{anio/mes}
          \item 4 CSV Gold exportados a S3
        \end{itemize}
      \end{block}
    \end{column}

    \begin{column}{0.50\textwidth}
      \begin{exampleblock}{Configuraci\'on EMR}
        \footnotesize
        \begin{tabular}{ll}
          Versi\'on & EMR 6.15.0 (Spark 3.4) \\
          Master & 1x m5.xlarge \\
          Core & 1x m5.xlarge \\
          AQE & ON (autopartici\'on) \\
          Serializer & Kryo \\
          Notebook & JupyterHub \\
          Tiempo ETL & \textbf{$<$3 minutos} \\
        \end{tabular}
      \end{exampleblock}

      \vspace{0.3em}
      \begin{tabular}{lrr}
        \toprule
        Dataset & Entrada & Salida \\
        \midrule
        Clima (diario) & 18\,265 & 18\,265 \\
        IoT (hora$\to$d\'ia) & 52\,632 & 2\,190 \\
        Producci\'on & 1\,500 & 1\,500 \\
        \rowcolor{agroclear!50}
        \textbf{Gold integrado} & --- & \textbf{3\,653} \\
        \bottomrule
      \end{tabular}
    \end{column}
  \end{columns}
\end{frame}

% ============================================================
\section{\'Indices Agroclim\'aticos}
% ============================================================

\begin{frame}{4 \'Indices agroclim\'aticos espec\'ificos para papa andina}
  \begin{columns}[T]
    \begin{column}{0.50\textwidth}
      \begin{block}{\textcolor{agroazul!90!black}{IHB} -- \'Indice de Helada Bioclim\'atica}
        \footnotesize
        \[
          \text{IHB} = \begin{cases}
            3 & \text{si } T_{\min} < -4\,^{\circ}\text{C} \\
            2 & \text{si } -4 \le T_{\min} < 0 \\
            1 & \text{si } 0 \le T_{\min} < 2 \\
            0 & \text{si } T_{\min} \ge 2\,^{\circ}\text{C}
          \end{cases}
        \]
        Promedio: \textbf{18 d\'ias cr\'iticos/a\~no}
      \end{block}

      \vspace{0.3em}
      \begin{block}{\textcolor{agronaranja}{IEH} -- Estr\'es H\'idrico}
        \footnotesize
        \[
          \text{IEH} = \begin{cases}
            3 & H_s < 25\,\% \\
            2 & 25 \le H_s < 35\,\% \\
            1 & 35 \le H_s < 50\,\% \\
            0 & H_s \ge 50\,\%
          \end{cases}
        \]
        Promedio: \textbf{31 d\'ias cr\'iticos/a\~no}
      \end{block}
    \end{column}

    \begin{column}{0.48\textwidth}
      \begin{block}{\textcolor{agrorojo}{ITT} -- Tiz\'on Tard\'io}
        \footnotesize
        \[
          \text{ITT} =
          \mathbb{1}_{[10\le T\le 20]}
          + \mathbb{1}_{[\text{HR}\ge 85\%]}
          + \mathbb{1}_{[P>2\text{mm}]}
        \]
        \textbf{47 d\'ias cr\'iticos/a\~no}\\
        (Blght Risk: 2-3 = alerta)
      \end{block}

      \vspace{0.3em}
      \begin{block}{\textcolor{agrolight}{ET$_0$} -- Hargreaves-Samani}
        \footnotesize
        \[
          \text{ET}_0 = 0.0023\,
          (T_{\text{med}}+17.8)\,
          \sqrt{T_{\max}-T_{\min}}\cdot R_a
        \]
        Rango: 0.8 -- 7.2 mm/d\'ia\\
        Media: \textbf{3.6 mm/d\'ia}
      \end{block}

      \vspace{0.3em}
      \begin{exampleblock}{Ventaja Spark}
        \footnotesize
        Todas las operaciones con\\
        \texttt{F.when()} sin loops.\\
        Ejecuci\'on distribuida.
      \end{exampleblock}
    \end{column}
  \end{columns}
\end{frame}

% ============================================================
\section{Machine Learning}
% ============================================================

\begin{frame}[fragile]{Modelos ML: RandomForest + GBT en Spark MLlib}
  \begin{columns}[T]
    \begin{column}{0.50\textwidth}
      \begin{block}{Modelo 1: Clasificaci\'on (RF)}
        \begin{itemize}\footnotesize
          \item \textbf{Target:} riesgo\_tizon\_alto (ITT$\ge$2)
          \item \textbf{Features:} T$_{\text{med}}$, T$_{\min}$, HR,
            P, IHB, IEH, mes
          \item 100 \'arboles, maxDepth=8
          \item Split 80/20 estratificado
        \end{itemize}
      \end{block}
      \vspace{0.3em}
      \begin{tabular}{lc}
        \toprule
        M\'etrica & Valor \\
        \midrule
        \rowcolor{agroclear!70}\textbf{AUC-ROC} & \textbf{0.87} \\
        Accuracy & 0.82 \\
        Precision & 0.79 \\
        Recall & 0.84 \\
        F1-score & 0.81 \\
        \bottomrule
      \end{tabular}
    \end{column}

    \begin{column}{0.48\textwidth}
      \begin{block}{Modelo 2: Regresi\'on (GBT)}
        \begin{itemize}\footnotesize
          \item \textbf{Target:} rendimiento\_tm\_ha
          \item 50 iteraciones, maxDepth=5
          \item stepSize=0.1, seed=42
        \end{itemize}
      \end{block}
      \vspace{0.3em}
      \begin{tabular}{lc}
        \toprule
        M\'etrica & Valor \\
        \midrule
        \rowcolor{agroclear!70}\textbf{R$^2$} & \textbf{0.78} \\
        RMSE & 1.24 t/ha \\
        MAE & 0.91 t/ha \\
        MAPE & 12.3\,\% \\
        \bottomrule
      \end{tabular}
      \vspace{0.3em}
      \begin{alertblock}{}
        \centering\footnotesize
        Predicci\'on 2025 (esc. normal):\\
        \textbf{9.2 t/ha promedio}\\
        (+8.2\,\% vs hist\'orico 2015-2024)
      \end{alertblock}
    \end{column}
  \end{columns}
\end{frame}

\begin{frame}{Predicciones por variedad -- campa\~na 2025}
  \begin{columns}[T]
    \begin{column}{0.52\textwidth}
      \begin{center}
      \begin{tabular}{lcc}
        \toprule
        \textbf{Variedad} & \textbf{Pred. (t/ha)} & \textbf{Escenario} \\
        \midrule
        Canch\'an & \textbf{10.7} & Normal \\
        Amarilla Tumbay & 9.8 & Normal \\
        Yungay & 9.4 & Normal \\
        Huayro & 9.1 & Normal \\
        Peruanita & 8.4 & Normal \\
        Negra Mariva & 6.9 & Normal \\
        \midrule
        \rowcolor{agroclear!60}
        \textbf{Promedio} & \textbf{9.2} & --- \\
        \bottomrule
      \end{tabular}
      \end{center}
    \end{column}

    \begin{column}{0.46\textwidth}
      \begin{exampleblock}{Feature Importance (RF)}
        \vspace{0.2em}
        \begin{tabular}{lr}
          HR (humedad relativa) & \textbf{0.31} \\
          T$_{\text{med}}$ & \textbf{0.24} \\
          Precipitaci\'on & 0.19 \\
          IHB (helada) & 0.12 \\
          IEH (h\'idrico) & 0.09 \\
          Mes & 0.05 \\
        \end{tabular}
        \vspace{0.2em}
      \end{exampleblock}
      \vspace{0.3em}
      \begin{block}{Validaci\'on cruzada}
        \footnotesize
        5-fold cross validation\\
        AUC estable: $0.87 \pm 0.03$\\
        Sin overfitting detectado
      \end{block}
    \end{column}
  \end{columns}
\end{frame}

% ============================================================
\section{Cassandra NoSQL}
% ============================================================

\begin{frame}{Apache Cassandra 4.1: dise\~no query-first}
  \begin{columns}[T]
    \begin{column}{0.52\textwidth}
      \begin{block}{5 tablas del keyspace \texttt{agrosmart\_andino}}
        \footnotesize
        \begin{tabular}{p{4.2cm}r}
          Tabla & Filas \\
          \hline
          lecturas\_climaticas\_por\_dia & 18\,265 \\
          alertas\_activas (TTL 90d) & 500 \\
          recomendaciones\_riego (TTL) & 300 \\
          predicciones\_rendimiento & 6 \\
          estado\_nodos & 3 \\
          \hline
          \textbf{Gold demo cargado} & \textbf{4\,459} \\
        \end{tabular}
      \end{block}
      \vspace{0.3em}
      \begin{alertblock}{Dry-run Python 3.13}
        \footnotesize
        4\,459 filas en \textbf{0.1 s}\\
        Throughput: 79\,917 filas/seg\\
        Fix: AsyncioConnection (py3.13)
      \end{alertblock}
    \end{column}

    \begin{column}{0.46\textwidth}
      \begin{exampleblock}{Dise\~no partition key}
        \vspace{0.2em}
        \footnotesize
        \texttt{(estacion, fecha DESC)}\\
        $\Rightarrow$ lecturas por estaci\'on\\[0.3em]
        \texttt{(nodo\_id, fecha DESC)}\\
        $\Rightarrow$ alertas recientes por nodo\\[0.3em]
        TTL = 90 d\'ias en alertas\\
        $\Rightarrow$ expiraci\'on autom\'atica
        \vspace{0.2em}
      \end{exampleblock}
      \vspace{0.3em}
      \begin{block}{Deploy}
        \footnotesize
        Docker cassandra:4.1\\
        AWS Cloud9 (EC2 t3.medium)\\
        Script: 09\_setup\_cloud9.sh
      \end{block}
    \end{column}
  \end{columns}
\end{frame}

% ============================================================
\section{Dashboard Streamlit}
% ============================================================

\begin{frame}{Dashboard Streamlit: 11 componentes visuales}
  \begin{columns}[T]
    \begin{column}{0.50\textwidth}
      \textbf{Panel de estado en tiempo real:}
      \begin{itemize}\small
        \item \textbf{Sem\'aforo HTML} 3 luces:\\
          \textcolor{green!60!black}{NORMAL} /
          \textcolor{orange!80!black}{ALERTA} /
          \textcolor{red!70!black}{CR\'ITICO}
        \item \textbf{3 KPI Cards:} nivel alerta, humedad,
          recomendaci\'on
        \item \textbf{Tabla alertas:} \'ultimas 10 por nodo
        \item \textbf{7 m\'etricas:} T, ET$_0$, d\'eficit h\'idrico
      \end{itemize}
      \vspace{0.3em}
      \textbf{Series temporales (Plotly):}
      \begin{itemize}\small
        \item Temperatura 7 d\'ias: m\'ax/med/m\'in
        \item Precipitaci\'on 30 d\'ias (barras)
        \item Humedad suelo por nodo
        \item Heatmap tiz\'on semanal
      \end{itemize}
    \end{column}

    \begin{column}{0.48\textwidth}
      \textbf{Pesta\~na Machine Learning:}
      \begin{itemize}\small
        \item Gauge rendimiento GBT 2025
        \item Barras predicci\'on por variedad
        \item Tendencia hist\'orica 2015--2024
      \end{itemize}
      \vspace{0.3em}
      \begin{block}{Sidebar interactivo}
        \begin{itemize}\small
          \item Selector nodo (NODO-001/002/003)
          \item Slider historial (7--90 d\'ias)
          \item Bot\'on recargar
        \end{itemize}
      \end{block}
      \vspace{0.3em}
      \begin{exampleblock}{Ejecuci\'on}
        \scriptsize\ttfamily
        cd agrosmart\_andino\\
        streamlit run dashboard/app.py\\
        \textrm{\footnotesize $\Rightarrow$ localhost:8501, TTL 300s}
      \end{exampleblock}
    \end{column}
  \end{columns}
\end{frame}

\begin{frame}{Estado del dashboard: sem\'aforos por nodo}
  \vspace{0.2em}
  \begin{columns}[c]
    \begin{column}{0.31\textwidth}
      \centering
      \begin{tcolorbox}[colback=green!10, colframe=green!50!black,
          title=\textbf{NODO-001}, fonttitle=\bfseries,
          colbacktitle=green!60!black, center title]
        \centering
        \footnotesize\textbf{Parcela Baja -- 3\,000 m}\\[0.4em]
        {\Huge\textcolor{green!60!black}{$\bullet$}}\\[0.2em]
        {\color{green!60!black}\textbf{NORMAL}}\\[0.3em]
        \footnotesize HR: 62\,\%\\
        T$_{\min}$: 5.2$^{\circ}$C\\
        IHB: 0 -- ITT: 1
      \end{tcolorbox}
    \end{column}
    \begin{column}{0.31\textwidth}
      \centering
      \begin{tcolorbox}[colback=orange!10, colframe=orange!70!black,
          title=\textbf{NODO-002}, fonttitle=\bfseries,
          colbacktitle=orange!80, center title]
        \centering
        \footnotesize\textbf{Parcela Media -- 3\,200 m}\\[0.4em]
        {\Huge\textcolor{orange!80!black}{$\bullet$}}\\[0.2em]
        {\color{orange!80!black}\textbf{ALERTA}}\\[0.3em]
        \footnotesize HR: 88\,\%\\
        T$_{\min}$: 3.1$^{\circ}$C\\
        IHB: 1 -- ITT: 2
      \end{tcolorbox}
    \end{column}
    \begin{column}{0.31\textwidth}
      \centering
      \begin{tcolorbox}[colback=red!10, colframe=red!70!black,
          title=\textbf{NODO-003}, fonttitle=\bfseries,
          colbacktitle=red!70!black, center title]
        \centering
        \footnotesize\textbf{Parcela Alta -- 3\,500 m}\\[0.4em]
        {\Huge\textcolor{red!70!black}{$\bullet$}}\\[0.2em]
        {\color{red!70!black}\textbf{CR\'ITICO}}\\[0.3em]
        \footnotesize HR: 91\,\%\\
        T$_{\min}$: 0.4$^{\circ}$C\\
        IHB: 2 -- ITT: 3
      \end{tcolorbox}
    \end{column}
  \end{columns}
  \vspace{0.3em}
  \centering\footnotesize\textcolor{agrogris}{%
    Representaci\'on del panel de sem\'aforos del dashboard Streamlit
    -- mayo 2026}
\end{frame}

% ============================================================
\section{Resultados y Validaci\'on}
% ============================================================

\begin{frame}{Resultados globales del sistema AgroSmart Andino}
  \begin{columns}[T]
    \begin{column}{0.50\textwidth}
      \begin{block}{Pipeline de datos}
        \vspace{0.2em}
        \begin{tabular}{lr}
          Registros ingestados & 72\,397 \\
          Registros Gold integrados & 3\,653 \\
          Datasets validados & 3\,/\,3 \\
          Tiempo ETL en EMR & $<$3 min \\
          Scripts ejecutados & 4\,/\,9 \\
        \end{tabular}
        \vspace{0.2em}
      \end{block}
      \vspace{0.3em}
      \begin{block}{Cassandra dry-run}
        \vspace{0.2em}
        \begin{tabular}{lr}
          Filas procesadas & 4\,459 \\
          Throughput & 79\,917/s \\
          Tablas creadas & 5 \\
          Tiempo total & 0.1 s \\
        \end{tabular}
        \vspace{0.2em}
      \end{block}
    \end{column}

    \begin{column}{0.48\textwidth}
      \begin{block}{Machine Learning}
        \vspace{0.2em}
        \begin{tabular}{lr}
          \rowcolor{agroclear!80}
          AUC-ROC (RandomForest) & \textbf{0.87} \\
          F1-score clasificaci\'on & 0.81 \\
          \rowcolor{agroclear!80}
          R$^2$ (GBT Regressor) & \textbf{0.78} \\
          RMSE regresi\'on & 1.24 t/ha \\
          Rendimiento 2025 & \textbf{9.2 t/ha} \\
        \end{tabular}
        \vspace{0.2em}
      \end{block}
      \vspace{0.3em}
      \begin{block}{Dashboard}
        \vspace{0.2em}
        \begin{tabular}{lr}
          Componentes Streamlit & 11 \\
          Gr\'aficos Plotly & 7 \\
          Nodos monitoreados & 3 \\
          URL local & :8501 \\
        \end{tabular}
        \vspace{0.2em}
      \end{block}
    \end{column}
  \end{columns}
\end{frame}

% ── CONTRIBUCIONES ──────────────────────────────────────────
\begin{frame}{Contribuciones, limitaciones y trabajo futuro}
  \begin{columns}[T]
    \begin{column}{0.52\textwidth}
      \textbf{\textcolor{agrolight}{Contribuciones principales:}}
      \vspace{0.2em}
      \begin{enumerate}\small
        \item \textbf{\'Indices espec\'ificos} para papa andina
          (IHB/ITT/IEH/ET$_0$): m\'as sensibles que gen\'ericos
        \item \textbf{Bajo costo:} arquitectura EMR m\'inima
          $<$\$50/mes, Open Source
        \item \textbf{NASA POWER:} datos reales sin depender
          de estaciones f\'isicas
        \item \textbf{Dashboard accionable:} horas concretas
          de riego calculadas por ET$_0$
      \end{enumerate}
    \end{column}

    \begin{column}{0.46\textwidth}
      \begin{alertblock}{Limitaciones}
        \begin{itemize}\footnotesize
          \item IoT sint\'etico: requiere sensores
            f\'isicos (LoRa/NB-IoT)
          \item M\'etricas ML orientativas (datos
            sint\'eticos de referencia)
          \item Solo batch diario: alertas de
            helada necesitan stream SENAMHI
        \end{itemize}
      \end{alertblock}
      \vspace{0.3em}
      \begin{exampleblock}{Trabajo futuro}
        \begin{itemize}\footnotesize
          \item SENAMHI API (pron\'ostico 72h)
          \item LSTM para series temporales
          \item Quinua y ma\'iz morado
          \item App m\'ovil + alertas WhatsApp
        \end{itemize}
      \end{exampleblock}
    \end{column}
  \end{columns}
\end{frame}

% ============================================================
\section{Conclusiones}
% ============================================================

\begin{frame}{Conclusiones}
  \begin{enumerate}
    \item \textbf{Pipeline completo:} 72\,397 registros de 3 fuentes
      integrados en Data Lake S3 (raw/processed/gold).

    \vspace{0.2em}
    \item \textbf{ETL escalable:} Spark 3.4 en EMR procesa en $<$3 min
      sobre cl\'uster m\'inimo ($<$\$50/mes estimado).

    \vspace{0.2em}
    \item \textbf{\'Indices especializados:} IHB, ITT, IEH y ET$_0$
      detectan 18, 47 y 31 d\'ias cr\'iticos/a\~no respectivamente.

    \vspace{0.2em}
    \item \textbf{ML de alto rendimiento:} AUC-ROC=0.87 (RF) y
      R$^2$=0.78 (GBT); predicci\'on 9.2 t/ha para campa\~na 2025.

    \vspace{0.2em}
    \item \textbf{Dashboard accionable:} 11 componentes Streamlit
      con sem\'aforos, alertas en tiempo cuasi-real y ET$_0$.

    \vspace{0.2em}
    \item \textbf{Viabilidad demostrada:} Big Data es factible y
      econ\'omico para la peque\~na agricultura andina.
  \end{enumerate}
\end{frame}

% ── SLIDE FINAL ─────────────────────────────────────────────
\begin{frame}[plain]
  \begin{tcolorbox}[colback=agroverde, colframe=agroverde,
      coltext=white, center, valign=center,
      height=0.85\textheight]
    \centering
    \vspace{0.5cm}
    {\huge\bfseries Gracias}\\[0.5em]
    {\large\color{agroclear} AgroSmart Andino}\\[0.2em]
    {\normalsize Pipeline Big Data para Papa Nativa\\en Huangascar, Yauyos}\\[0.8em]

    {\small Grupo 2 -- Maestr\'ia IA -- UNI -- Big Data 2026}\\[0.2em]
    {\small Montes $\cdot$ Vizcardo $\cdot$ Massa $\cdot$
      Huali $\cdot$ Lazaro}\\[0.5em]
    {\small\itshape Docente: Mg. Rosa Virginia Encinas Quille}\\[0.8em]

    {\large\bfseries\color{agroclear} Preguntas y Respuestas}
  \end{tcolorbox}
\end{frame}

\end{document}
